\section{Energy normalization and the common coupling}\label{sec:inputs}

The proof needs more than the metric-measure limit: it needs convergence of
inverse operators and uniform recovery of spatial test functions with the
correct energy.  We state these inputs from the spectral companion
\cite{SC} precisely, keeping its unhalved electrical convention separate
from the Brownian convention of the theorem.

\subsection{The graph energy and the spatial scale}

Let \(P_n\) be the one-step transition operator of the walk, let
\(L_n=P_n-I\), and put
\[
 G_n=-nL_n=n(I-P_n).
\]
Thus \(-G_n\) is the generator with total attempt rate \(n\), including
attempts along loops.  It is self-adjoint on \(L^2(\mu_n)\).  For a real
vertex function \(u\), write
\(\E_n(u)=\ip{u}{G_nu}_{L^2(\mu_n)}\), and use polarization for the
bilinear form.  Choose either orientation of each edge and write
\(d_eu=u(e^+)-u(e^-)\).

\begin{lemma}[Graph energy]\label{lem:graph-energy}
For all real vertex functions \(u,v\),
\begin{equation}\label{eq:graph-energy}
 \E_n(u,v)=\frac12\sum_{e\in E(M_n)}d_eu\,d_ev.
\end{equation}
Every parallel edge is a separate summand, every loop has zero summand,
and the kernel of \(G_n\) consists of the constants.
\end{lemma}

\begin{proof}
Since \(\mu_n(x)=\deg(x)/(2n)\), every incident half-edge has
\(\mu_n\)-weighted accelerated rate \(1/2\).  A nonloop edge with
endpoints \(x,y\) therefore contributes
\[
 \frac12\bigl(u(x)(v(x)-v(y))+u(y)(v(y)-v(x))\bigr)
 =\frac12d_eu\,d_ev
\]
to \(\ip{u}{G_nv}_{L^2(\mu_n)}\).  Both half-edges of a loop contribute
zero.  Summing proves \eqref{eq:graph-energy}, including multiplicities.
Zero energy forces equality across each edge, hence constancy by
connectedness.
\end{proof}

On the limiting conformal sphere, write \(\E_h(f)=\E_h(f,f)\), also for
the closure, and let \(G\) be its nonnegative self-adjoint operator.  The
existence and identification of this closure are part of
\cite[Proposition~8.3]{SC}.  In the companion's notation the unhalved
operators are
\begin{equation}\label{eq:raw-operators}
 B_n=2G_n=-2nL_n,\qquad B_h=2G.
\end{equation}
In particular, the inverses imported below are inverses of \(2G_n\) and
\(2G\), not of \(G_n\) and \(G\).

The companion fixes space by a deterministic diameter normalization.  The
next elementary comparison, reproduced from \cite[Lemma~2.3]{SC}, shows
why its conclusions apply to the sequence \(a_n\) in the theorem.
For a finite real random variable \(Y\), let \(\mathfrak m(Y)\) be the
unique zero of \(r\mapsto\EE[\arctan(Y-r)]\).  Set
\[
 \Delta_n=\max\{1,\operatorname{diam}(V_n,d_n)\},\qquad
 \Delta=\operatorname{diam}(S,D_h),
\]
and define
\begin{equation}\label{eq:diameter-normalization}
 \widehat a_n=
 \exp\{\mathfrak m(\log\Delta)-\mathfrak m(\log\Delta_n)\}.
\end{equation}
The expectations defining these centers use the unconditional laws, so
\(\widehat a_n\) is deterministic.

\begin{lemma}[Comparison of spatial normalizations]
\label{lem:space-normalization}
Under the metric-measure convergence assumption of the theorem,
\(\widehat a_n/a_n\to1\).
\end{lemma}

\begin{proof}
For every finite real \(Y\), bounded convergence shows that
\(F_Y(r)=\EE[\arctan(Y-r)]\) is continuous, strictly decreasing, and
has limits \(\pi/2\) and \(-\pi/2\) at the two ends of the real line.
Its zero therefore exists and is unique.  Directly,
\(\mathfrak m(Y+c)=\mathfrak m(Y)+c\).  If \(Y_j\) converges in law
to \(Y\), then \(F_{Y_j}(r)\to F_Y(r)\) for each fixed \(r\).
For each \(\varepsilon>0\), the values at
\(\mathfrak m(Y)-\varepsilon\) and
\(\mathfrak m(Y)+\varepsilon\) have opposite strict signs.  The same
eventually holds for \(F_{Y_j}\), trapping its zero between these two
points.  Thus \(\mathfrak m\) is continuous under convergence in law.

Diameter is continuous in the Gromov--Hausdorff topology.  Since
\(a_n\to0\), the assumed convergence gives
\[
 a_n\Delta_n\ \Longrightarrow\ \Delta,
 \qquad 0<\Delta<\infty\quad\hbox{almost surely}.
\]
Applying the preceding continuity statement after taking logarithms and
using the shift identity gives
\[
 \log\frac{a_n}{\widehat a_n}
 =\mathfrak m\bigl(\log(a_n\Delta_n)\bigr)
      -\mathfrak m(\log\Delta)\longrightarrow0.
\]
\end{proof}

\subsection{The geometric and analytic inputs}

We use \(\Sph\) for the coordinate sphere and \(\rho\) for its round
distance.  For each occurrence of a vertex, an incident edge, and an
incident face, form an equilateral triangle with vertices labelled by
these three objects.  Repeated incidences remain distinct.  Gluing the
\(4n\) triangles according to the cyclic incidences gives a topological
sphere with its piecewise Euclidean conformal structure.  Uniformize this
surface onto \(\Sph\), with the normalization of \cite{SC}, and write
\(z_n:V_n\to\Sph\) for the primal vertex locations.  The images of the
triangles are the embedded \emph{flags}; the primal edges are embedded
along their vertex-to-edge-midpoint sides.  This common coordinate will
retain both the intrinsic metric and the conformal energy.

\begin{proposition}[Inputs from the spectral companion]
\label{prop:inputs}
From every subsequence of sizes tending to infinity there is a further
subsequence and a representation coupling on which the following
properties hold almost surely, simultaneously.  The limiting sphere in
these coordinates has the ordinary unit-area \(\sqrt3\)-LQG law, with
area \(\mu_h\) and metric \(D_h\); the latter induces the round topology.

\begin{enumerate}[label=\textup{(\roman*)}]
\item \emph{Geometry in one coordinate.}  The vertex measures satisfy
\((z_n)_*\mu_n\to\mu_h\) weakly, and
\begin{align}
 \delta_n&:=\sup_{u,v\in V_n}
  \abs{a_nd_n(u,v)-D_h(z_n(u),z_n(v))}\longrightarrow0,
       \label{eq:vertex-distortion}\\
 \varepsilon_n&:=d_{\mathrm H}^{D_h}(z_n(V_n),\Sph)
       \longrightarrow0.\label{eq:vertex-density}
\end{align}
The maximum round diameter of the embedded flags tends to zero.
Consequently the maximum round diameter of an embedded primal edge tends
to zero as well.

\item \emph{One Hilbert space.}  There is an order
\(v_n(i)\), \(0\le i<2n\), of the primal corner occurrences and a
continuous tour \(\eta:[0,1]\to\Sph\) such that
\begin{equation}\label{eq:tour-tracking}
 \eta_*\dd t=\mu_h,\qquad
 \sup_{0\le t<1}
 \rho\bigl(z_n(v_n(\lfloor2nt\rfloor)),\eta(t)\bigr)
 \longrightarrow0.
\end{equation}
Each vertex \(v\) occurs \(\deg(v)\) times in the corner list.  For
\(\mu_h\)-almost every \(z\), the fiber \(\eta^{-1}(\{z\})\)
contains exactly one time.  With \(H=L^2([0,1],\dd t)\), the maps
\begin{align*}
 J_n:L^2(\mu_n)&\longrightarrow H,&
 (J_nu)(t)&=u(v_n(\lfloor2nt\rfloor)),\quad 0\le t<1,\\
 J:L^2(\mu_h)&\longrightarrow H,& Jf&=f\circ\eta
\end{align*}
are isometries, and \(J\) is unitary.  In particular,
\begin{equation}\label{eq:continuous-pullbacks}
 \norm{J_n(f\circ z_n)-Jf}_{L^\infty([0,1])}\longrightarrow0
 \qquad(f\in C(\Sph)).
\end{equation}

\item \emph{Centered inverse convergence.}  The closed continuum form
\(\E_h\) is a regular Dirichlet form and is the form of Liouville
Brownian motion with the stated half-energy clock.  Its operator \(G\)
has kernel exactly the constants.  For a nonnegative operator whose
kernel is the constants, write \(A^\dagger\) for its inverse on the
mean-zero subspace, extended by zero on constants.  Let \(\Pi\) be the
orthogonal projection onto the constants in \(H\), and put
\begin{equation}\label{eq:centered-inverses}
 K_n=J_n(2G_n)^\dagger J_n^*,\qquad
 K=J(2G)^\dagger J^*.
\end{equation}
Thus \(K_n\) is also zero on \(J_n(L^2(\mu_n))^\perp\).  These are
bounded positive self-adjoint operators, and
\begin{equation}\label{eq:inverse-convergence}
 \norm{K_n-K}_{\mathrm{HS}}\longrightarrow0.
\end{equation}

\item \emph{Local recovery with unit conductivity.}  There are ordinary
conformal chart domains whose interiors cover \(\Sph\).  Let
\(V\Subset W\Subset U\) be open sets with \(U\) compactly contained
in one of these charts, and let \(\chi\) be a real smooth function on
\(\Sph\) satisfying \(0\le\chi\le1\) and
\(\operatorname{supp}\chi\subset V\).  There are local vertex
functions \(\widetilde w_n\), defined on graphs containing all
incidences needed over \(W\), that converge uniformly to \(\chi\) on
a neighborhood of \(\overline W\), in the vertex locations, and satisfy
\begin{equation}\label{eq:local-recovery}
 \limsup_n\sum_{e:\,z_n(e)\in W}
       (d_e\widetilde w_n)^2
 \le\int_U\abs{\nabla\chi}^2\,\dd\mathrm{vol}.
\end{equation}
Here \(z_n(e)\) is any chosen point on the embedded edge, and the
integral is the unhalved conformal Dirichlet integral.  Every edge is
counted with its multiplicity.  If one endpoint lies in a fixed compact
subset of \(W\), then eventually its entire edge, both endpoints, and
all the required incidences lie in \(W\), uniformly over such edges.
\end{enumerate}
The assertions in \textup{(iv)} hold simultaneously for these buffered
tests on the coupled outcome, so the tests may be chosen after fixing
that outcome.
\end{proposition}

\begin{proof}[Source of the inputs and comparison of conventions]
The substantial assertions are imported from \cite{SC}.  Its
Proposition~2.4 supplies the common coordinate, the full flag mesh, and
the tours; its Theorem~5.2 supplies the local recovery assertion; its
Proposition~8.3 identifies the closed continuum form and centered inverse;
and its Lemma~9.1 and Equation~(9.11) give the Hilbert-space identification
and Hilbert--Schmidt convergence.  In the common extraction of its
Section~9, retain also the countable family of local constrained-energy
infima from its Lemma~5.1.  These are infima of finite sums of squared
edge differences under rational interval constraints on vertex values,
viewed in the compact space \([0,\infty]\).  Their joint limits supply
the simultaneous recovery statement of Theorem~5.2, and persist on
further subsequences.  Proposition~2.4 preserves the geometric
conclusions under this joint extraction, so all four assertions hold on
one coupling.

The model conventions agree exactly: the companion uses the same rooted
FK weight \(2^{\ell/2}\), counts all \(n\) primal edges, and pushes
normalized corner counting to \(\deg(v)/(2n)\).  Its all-edge network
therefore has operator \(B_n=2G_n\) by
Lemma~\ref{lem:graph-energy}.  Its limiting surface is the ordinary
unit-area sphere with the original marks forgotten.  Auxiliary marks and
rerootings used in the coupling change neither of these marginals.

For completeness, the companion first gives Hausdorff convergence of
the distance graphs
\[
 \{(z_n(u),z_n(v),\widehat a_nd_n(u,v)):u,v\in V_n\}
 \longrightarrow
 \{(x,y,D_h(x,y)):x,y\in\Sph\}
\]
in the product of the round topology and the ordinary topology on
\([0,\infty)\).  Uniform continuity of \(D_h\) on the compact
\(\Sph\times\Sph\) turns this into vanishing vertex distortion.
Approximating the triples \((x,x,0)\) in the opposite direction gives
round density of the vertex locations, hence \(D_h\)-density.
Lemma~\ref{lem:space-normalization} permits replacement of
\(\widehat a_n\) by \(a_n\): along this coupling the
\(\widehat a_n\)-scaled diameters are bounded, and the ratio of the
two scales tends to one.

The corner multiplicities give the isometry of \(J_n\), and the area
pushforward gives that of \(J\).  The almost-everywhere singleton-fiber
property makes \(J\) onto, as detailed in \cite[Lemma~9.1]{SC}.
Uniform tour tracking and uniform continuity of \(f\) give
\eqref{eq:continuous-pullbacks}.  The centered inverses in the companion
are those of the unhalved operators, which yields exactly
\eqref{eq:centered-inverses}.

Finally, the chart families in \cite[Proposition~2.4(iv)]{SC} use two
distinct exploration roots and separate auxiliary mark pairs.  Their
interiors cover the entire sphere, not merely the complement of one root.
The finite-sphere clause of its Theorem~5.2 gives
\eqref{eq:local-recovery} in the common coordinate, simultaneously for
continuous finite-energy inputs and buffered domains.  We have stated
only the smooth supported case needed here.  The vanishing flag mesh and
the strict buffers ensure that edges incident at the compact support are
indeed among the edges counted in \(W\).  Conformal and anticonformal
coordinate changes preserve the Dirichlet integral, and hence the
diffusion and its clock.
\end{proof}

Fix henceforth one outcome in the probability-one event of
Proposition~\ref{prop:inputs}, and relabel its further subsequence by
\(n\).  Until we return to random environments, all graphs, coordinates,
and the limiting measured conformal sphere are fixed.  The symbols
\(\PP\) and \(\EE\) will refer to the law and expectation of walks
sampled with fresh randomness on these fixed graphs, started from
\(\mu_n\), or of the limiting motion started from \(\mu_h\).
